\subsection{Properties of the Polynomial Algebra of an Open Set}
\label{sec:ch4-2-polynomial-algebra}

Quantum field theory provides a set of candidates for local measurements,
observables which correspond to field measurements performed in a laboratory
of finite size and completed in a finite time. These are the field operators
smeared with test functions of compact support. The present section is
concerned with the properties of an algebra associated with such quantities.
For simplicity, we again discuss the theory of a hermitian scalar field.

Let $\PolynomialAlgebra(\mathcal O)$ be the set of all polynomials of the form
\begin{equation}
  c+\sum_{j=1}^{N}
  \varphi\bigl(f_1^{(j)}\bigr)\cdots
  \varphi\bigl(f_j^{(j)}\bigr),
  \tag{4-9}\label{eq:ch4-9}
\end{equation}
where the $f_k^{(j)}$ are test functions whose support is in the open set
$\mathcal O$ of space-time, and $c$ is any complex constant. Clearly, if $p$
and $q$ are two such polynomials so are $p+q$, $\alpha p$, $p^\dagger$, and
$pq$. This means that $\PolynomialAlgebra(\mathcal O)$ is a $*$-algebra, the
\emph{polynomial algebra} of $\mathcal O$. The next theorem is a surprise:
the algebra associated with \emph{any} open set $\mathcal O$ has
$\ket{\Omega}$ as a cyclic vector. It is due to Reeh and Schlieder.

\begin{theorem}[4-2]\label{thm:ch4-2-reeh-schlieder}
Suppose $\mathcal O$ is an open set of space-time. Then $\ket{\Omega}$ is a
cyclic vector for $\PolynomialAlgebra(\mathcal O)$, if it is a cyclic vector
for $\PolynomialAlgebra(\R^4)$. That is, vectors of the form
\begin{equation}
  \sum_{j=0}^{N}
  \varphi\bigl(f_1^{(j)}\bigr)\cdots
  \varphi\bigl(f_j^{(j)}\bigr)\ket{\Omega}
  \tag{4-10}\label{eq:ch4-10}
\end{equation}
with $\supp f_j^{(j)}\subset\mathcal O$ are dense in $\Hilbert$.

\begin{proof}
Let $\ket{\Psi}$ be orthogonal to all vectors of the form \eqref{eq:ch4-10}. We are
going to prove that $\ket{\Psi}$ is also orthogonal to all polynomials in the
smeared field with no conditions on the supports of the test functions, that
is, that $\ket{\Psi}$ is orthogonal to
$\PolynomialAlgebra(\R^4)\ket{\Omega}=D_0$. This will then imply that
$\ket{\Psi}=0$, because by assumption $D_0$ spans the Hilbert space $\Hilbert$.
The method of proof is another typical application of the principle of
analytic continuation.

The first step in the proof is to argue that the symbolic expression
\begin{equation*}
  \braket{\Psi}{\Psi_f}
  =\int \bra{\Psi}\varphi(x_1)\varphi(x_2)\cdots
  \varphi(x_n)\ket{\Omega}
  f(x_1,x_2,\ldots,x_n)\,\dd x_1\cdots \dd x_n
\end{equation*}
has a meaning for all test functions, $f$, in $\Schwartz$, and any
$\ket{\Psi}$ in $\Hilbert$. The proof goes just like that in the remarks
after \eqref{eq:pct-3-24}. One first notes that
$\bra{\Psi}\varphi(f_1)\cdots\varphi(f_n)\ket{\Omega}$ is a continuous
multilinear functional of the test functions $f_1\ldots f_n$ and applies the
Schwartz nuclear theorem to extend it to a tempered distribution in all the
variables together. Thus there is a tempered distribution $F$ defined by
\begin{equation*}
  F(-x_1,x_1-x_2,\ldots,x_{n-1}-x_n)
  \equiv
  \bra{\Psi}\varphi(x_1)\cdots\varphi(x_n)\ket{\Omega}.
\end{equation*}

Standard arguments\footnote{%
See Section~\ref{sec:2-6}.}%\allowbreak
 show that the Fourier transform of $F$ vanishes\newline unless each four-momentum
 variable lies in the physical spectrum.

Therefore there exists a holomorphic function $F$, holomorphic in the tube
$\Tube_n$ in the variables $(-x_1)-\ii\eta_1$,
$(x_1-x_2)-\ii\eta_2,\ldots,(x_{n-1}-x_n)-\ii\eta_n$, whose boundary value
as $\eta_1\cdots\eta_n\to0$ in $V_+$ is $F$. This boundary value vanishes
for $-x_1,x_1-x_2,\ldots,x_{n-1}-x_n$ in an open set determined by
$x_1\ldots x_n\in\mathcal O$. Thus, by Theorem \ref{thm:edge-of-wedge-zero-boundary} $F$ vanishes and
therefore its boundary value vanishes for all $x_1,\ldots,x_n$. This shows
$\ket{\Psi}$ is orthogonal to $D_0$.
\end{proof}
\end{theorem}

It is well known in the theory of algebras of bounded operators that a cyclic
vector for a $*$-algebra, $\PolynomialAlgebra$, is a separating vector for its
commutant. That means: if the set of vectors $\PolynomialAlgebra\ket{\Psi}$ is
dense, then the equality $T\ket{\Psi}=0$ for an operator $T$ commuting with
all operators of $\PolynomialAlgebra$ implies $T=0$. (For a proof see
\cite{ch4-4}.)
Theorem \ref{thm:ch4-3-separating} is an analogue of this statement for our case. It can be
interpreted as meaning that it is difficult to isolate a system described by
fields from outside effects. The algebra $\PolynomialAlgebra(\mathcal O)$
never contains any annihilation or creation operators if the open set
satisfies a certain boundedness condition. To define the boundedness
condition, consider the set of all points which are space-like with respect to
every point of $\mathcal O$. The interior of this set is $\mathcal O'$. Then
we assert

\begin{theorem}[4-3]\label{thm:ch4-3-separating}
If $\mathcal O$ is an open set for which $\mathcal O'$ is not empty, and
$T\in\PolynomialAlgebra(\mathcal O)$, then
\begin{equation}
  T\ket{\Omega}=0
  \tag{4-11}\label{eq:ch4-11}
\end{equation}
implies $T=0$.

\begin{proof}
Let $\ket{\Phi}$ be a vector in the domain of definition $D$ of $\varphi$.
Let $\ket{\Psi}=P'\ket{\Omega}$, where
$P'\in\PolynomialAlgebra(\mathcal O')$. Then for any
$T\in\PolynomialAlgebra(\mathcal O)$ satisfying \eqref{eq:ch4-11},
\begin{equation}
  \braket{\Psi}{T^\dagger\Phi}
  =\braket{T\Psi}{\Phi}
  =\braket{TP'\Omega}{\Phi}
  =\braket{P'T\Omega}{\Phi}=0.
  \tag{4-12}\label{eq:ch4-12}
\end{equation}
where in the last equality \eqref{eq:ch4-11} has been used. Since, by Theorem \ref{thm:ch4-2-reeh-schlieder},
vectors of the form $\PolynomialAlgebra(\mathcal O')\ket{\Omega}$ span
$\Hilbert$, we have $T^\dagger\ket{\Phi}=0$. This in turn implies
$T\ket{\Psi}=0$, for $\ket{\Psi}\in D$ since
$\braket{T\Psi}{\Phi}=\braket{\Psi}{T^\dagger\Phi}$ and $D$ is dense.
\end{proof}
\end{theorem}

\emph{Remarks:}

\begin{enumerate}
\item Theorem \ref{thm:ch4-3-separating} retains its validity if, instead of assuming that
commutators of fields vanish at space-like separations, we assume that
anti-commutators vanish. This will be used in the proof of Theorem \ref{thm:ch4-8}.

\item Any bounded open set $\mathcal O$ has the property that $\mathcal O'$ is
non-empty, and so the theorem applies.
\end{enumerate}

Theorem \ref{thm:ch4-4-adjoined-projection}, also due to Reeh and Schlieder, shows that
$\PolynomialAlgebra(\mathcal O)$ becomes irreducible when a single operator
is adjoined.

\begin{theorem}[4-4]\label{thm:ch4-4-adjoined-projection}
Let $E_0$ be the projection operator onto the vacuum state, which is supposed
cyclic with respect to the field. Then, for every open set $\mathcal O$, the
set of operators $\{E_0,\PolynomialAlgebra(\mathcal O)\}$ is irreducible.

\begin{proof}
Recall the definition of an irreducible set, Eq.~\eqref{eq:ch3-irreducibility}. Suppose that for some
bounded operator, $C$, and all $\ket{\Phi},\ket{\Psi}\in D_0$, we have
\begin{equation}
  \matrixel{\Phi}{C\varphi(f)}{\Psi}
  =\matrixel{\varphi(f)^\dagger\Phi}{C}{\Psi},
  \tag{4-13}\label{eq:ch4-13}
\end{equation}
where $\supp f\subset\mathcal O$, and suppose that
\begin{equation}
  CE_0=E_0C.
  \tag{4-14}\label{eq:ch4-14}
\end{equation}

Then \eqref{eq:ch4-13} holds, in particular, for a state $\ket{\Psi}$ of the form
\begin{equation}
  \ket{\Psi}=p\ket{\Omega},\qquad
  p\in\PolynomialAlgebra(\mathcal O)
  \tag{4-15}\label{eq:ch4-15}
\end{equation}
and
\begin{equation}
\begin{aligned}
  \bra{\Phi}C\ket{\Psi}
    &=\bra{\Phi}Cp\ket{\Omega}
     =\bra{p^\dagger\Phi}C\ket{\Omega}
     =\bra{p^\dagger\Phi}CE_0\ket{\Omega} \\
    &=\bra{p^\dagger\Phi}E_0C\ket{\Omega}
     =\braket{p^\dagger\Phi}{\Omega}
      \bra{\Omega}C\ket{\Omega},
\end{aligned}
\tag{4-16}\label{eq:ch4-16}
\end{equation}
where \eqref{eq:ch4-13}, \eqref{eq:ch4-14}, and the definition of $E_0$,
\begin{equation*}
  E_0\ket{\Psi}=\braket{\Omega}{\Psi}\ket{\Omega},
\end{equation*}
have been used. Clearly, \eqref{eq:ch4-16} says
\begin{equation*}
  \bra{\Phi}C\ket{\Psi}=c_0\braket{\Phi}{\Psi},
\end{equation*}
where
\begin{equation*}
  c_0=\bra{\Omega}C\ket{\Omega}.
\end{equation*}

Since a bounded operator is continuous and states $\ket{\Phi},\ket{\Psi}$ are
dense in $\Hilbert$, this implies $C=c_0$.
\end{proof}
\end{theorem}

Theorem \ref{thm:ch4-4-adjoined-projection} shows that for any $\mathcal O$ the adjunction of $E_0$ to
$\PolynomialAlgebra(\mathcal O)$ produces an irreducible set of operators.
The next theorem shows that for a special choice of $\mathcal O$, namely, the
whole of space-time, $E_0$ is effectively contained in
$\PolynomialAlgebra(\mathcal O)$, so $\PolynomialAlgebra(\mathcal O)$ itself
is irreducible. We give a proof for a theory of a neutral scalar field. The
extension to a set of fields of arbitrary transformation laws is easy.

\begin{theorem}[4-5]\label{thm:ch4-5-smeared-fields}
In any field theory, the smeared fields form an irreducible set of operators.
\end{theorem}

\emph{Remark:}

Recall that with our definition of a field theory the vacuum is a cyclic
vector for the smeared fields; this hypothesis is essential for the truth of
the theorem. The theorem provides the promised justification of the
definition of field theory in terms of cyclicity of the vacuum rather than
irreducibility of the smeared fields.

\begin{proof}
Notice that if $C$ is a bounded operator which satisfies
\begin{equation*}
  \matrixel{\Phi}{C\varphi(f)}{\Psi}
  =\matrixel{\varphi(f)^\dagger\Phi}{C}{\Psi}
\end{equation*}
for all $\ket{\Phi},\ket{\Psi}\in D_0$, then $C$ satisfies
\begin{equation}
  \matrixel{\Phi}{C\varphi(f_1)\cdots\varphi(f_n)}{\Psi}
  =\matrixel{\varphi(f_n)^\dagger\cdots\varphi(f_1)^\dagger\Phi}{C}{\Psi}.
  \tag{4-17}\label{eq:ch4-17}
\end{equation}

Consider, in particular, the case
\begin{equation}
\begin{aligned}
  &\bra{\Omega}C\varphi(\{a,1\}f_1)\cdots
  \varphi(\{a,1\}f_n)\ket{\Omega} \\
  &\quad=\bra{\varphi(\{a,1\}f_n)^\dagger\cdots
  \varphi(\{a,1\}f_1)^\dagger\Omega}C\ket{\Omega}.
\end{aligned}
\tag{4-18}\label{eq:ch4-18}
\end{equation}

Rewritten using the transformation law of $\varphi$ and the invariance of the
vacuum state under translation, it is
\begin{equation*}
  \bra{\Omega}CU(a,1)\varphi(f_1)\cdots\varphi(f_n)\ket{\Omega}
  =\bra{\varphi(f_n)^\dagger\cdots\varphi(f_1)^\dagger\Omega}
  U(-a,1)C\ket{\Omega}.
\end{equation*}

By Fourier-transforming this equation in $a$ we get the assertion
\begin{equation*}
  \bra{\Omega}CE(S)\varphi(f_1)\cdots\varphi(f_n)\ket{\Omega}
  =\bra{\varphi(f_n)^\dagger\cdots\varphi(f_1)^\dagger\Omega}
  E(-S)C\ket{\Omega},
\end{equation*}
where $S$ is any measurable set in momentum space. This is explained in
Section~\ref{sec:2-6}.

But since the spectrum of energy momentum lies in or on the plus cone, if
$-S$ lies in the physical spectrum and does not include $p=0$, the left-hand
side vanishes. That means that $C\ket{\Omega}$ is orthogonal to all states
$E(-S)\varphi(f_n)^\dagger\cdots\varphi(f_1)^\dagger\ket{\Omega}$ and therefore
$C\ket{\Omega}=c\ket{\Omega}$, where $c$ is a complex number. Then \eqref{eq:ch4-17} for
the special case $\ket{\Psi}=\ket{\Omega}$ is
\begin{equation*}
\begin{aligned}
  \matrixel{\Phi}{C\varphi(f_1)\cdots\varphi(f_n)}{\Omega}
  &=\matrixel{\varphi(f_n)^\dagger\cdots\varphi(f_1)^\dagger\Phi}{C}{\Omega} \\
  &=c\matrixel{\Phi}{\varphi(f_1)\cdots\varphi(f_n)}{\Omega}.
\end{aligned}
\end{equation*}

This implies
\begin{equation*}
  \bra{\Phi}C\ket{\Psi}=c\braket{\Phi}{\Psi}
\end{equation*}
for all $\ket{\Phi},\ket{\Psi}\in D_0$, whence by the continuity of $C$ and
the denseness of $D_0$, $C=c$.
\end{proof}
